\subsection{Retour sur la partie algorithmique}

Dans la partie algorithmique, nous avons rencontré plusieurs difficultés tout au long du projet.

Pour commencer, certains tests ont été réalisés, néanmoins il ne s’agissait pas de tests unitaires au sens strict. Les tests que nous avons faits avaient avant tout pour but de visualiser nos programmes au fur et à mesure du développement, comme par exemple la création d’un labyrinthe fictif pour vérifier la bonne génération des ordres pour un chemin donné. Bien que ces tests nous aient été utiles, la détection d’erreurs plus spécifiques est à revoir.

En ce qui concerne l’algorithme de résolution du labyrinthe, sa logique est correcte et l’algorithme est correct. Malheureusement, nous avons rencontré de nombreux problèmes au moment de l’implémentation, notamment concernant les allocations dynamiques de mémoire. Ce problème est avant tout lié au manque de maîtrise de certains concepts du langage C. Manquant de temps, nous n’avons malheureusement pas pu compléter cette partie.

Enfin, le problème principal que nous avons pu relever dans ce projet est le manque d’organisation et la gestion du temps. Nous avions identifié ces problématiques dès le départ, mais leur impact s’est surtout fait ressentir pendant les dernières semaines avant le rendu, quand les problèmes ont commencé à s’accumuler. Cette difficulté nous a montré l’importance d’une meilleure organisation pour les projets à venir.

Pour conclure, ces problèmes nous ont permis de nous rendre compte de plusieurs points. Ils ont mis en avant l’importance des tests unitaires, d’une meilleure planification, et des points à revoir en programmation.
